\section{Potencial de Aplicação em Veículos Hipersónicos: Análise Zonal}\label{sec:potencial_aplicacao_hipersonica}

O voo sustentado em regimes hipersônicos submete as diferentes zonas de uma aeronave a perfis aerotérmicos e mecânicos heterogêneos. Enquanto superligas convencionais à base de níquel sofrem amolecimento acentuado acima de $1000~^\circ\mathrm{C}$ e cerâmicas monolíticas falham por fragilidade mecânica e choque térmico, as ligas de alta entropia oferecem flexibilidade composicional e microestrutural customizável para cada subsistema do veículo.

Nas seções frontais de estagnação, representadas pelos bordos de ataque e cones de nariz, as temperaturas de choque ultrapassam com frequência $2000~^\circ\mathrm{C}$ sob fluxo aerodinâmico ionizado. Nesse domínio, RHEAs de matriz BCC como TiZrHfNbTa e NbMoTaW combinam altíssimo ponto de fusão ($> 2200~^\circ\mathrm{C}$) com tenacidade metálica suficiente para absorver impactos dinâmicos sem propagação catastrófica de trincas.

Na seção de propulsão \emph{scramjet} e canais de combustão interna, onde gases quentes exercem pressão mecânica e química severa em temperaturas contínuas de até $1800~^\circ\mathrm{C}$, o efeito cinético de difusão lenta intrínseco às HEAs atua como freio natural contra a deformação por fluência. Essa estabilidade supera em ordens de magnitude a integridade temporal das ligas metálicas tradicionais. Já nos painéis de fuselagem e superfícies de controle aerodinâmico, a prioridade reside em manter a rigidez estrutural sob dilatação diferencial, mantendo mais de 80\% da resistência à temperatura ambiente sem empenamento excessivo.

Por outro lado, o armazenamento de hidrogênio líquido em tanques criogênicos expõe componentes a temperaturas próximas de $-253~^\circ\mathrm{C}$ ($20~\mathrm{K}$). Ao contrário de aços ferríticos e refratários comuns que se tornam extremamente frágeis sob frio severo, formulações FCC de alta entropia ativam mecanismos sinérgicos de maclação e discordâncias, aumentando simultaneamente a resistência e a ductilidade no extremo criogênico. O mapeamento comparativo das zonas e de suas soluções metalúrgicas encontra-se consolidado na Tabela~\ref{tab:comparativo_zonal_hea}.

\begin{table}[!htbp]
\centering
\caption{Resumo das aplicações de Ligas de Alta Entropia por região funcional do veículo hipersônico.}\label{tab:comparativo_zonal_hea}
\small
\begin{tabularx}{\textwidth}{>{\raggedright\arraybackslash}p{3.2cm} >{\raggedright\arraybackslash}X >{\raggedright\arraybackslash}X}
\toprule
\textbf{Região do Veículo} & \textbf{Desafio Operacional} & \textbf{Vantagem Tecnológica da HEA/RHEA} \\
\midrule
\textbf{Bordos de Ataque e Nariz} & Choque aerodinâmico, atrito e calor de estagnação ($> 2000~^\circ\mathrm{C}$). & Alta refratariedade e retenção de tenacidade metálica contra choque térmico. \\
\textbf{Propulsão \emph{Scramjet}} & Gases corrosivos em alta pressão e temperaturas de até $1800~^\circ\mathrm{C}$. & Cinética atômica lenta bloqueando deformação plástica contínua por fluência. \\
\textbf{Fuselagem e Painéis} & Cargas dinâmicas de voo e gradientes térmicos superficiais. & Elevada rigidez específica com baixa dilatação térmica prejudicial. \\
\textbf{Tanques Criogênicos} & Contato com propelente líquido a $-253~^\circ\mathrm{C}$ ($20~\mathrm{K}$). & Mecanismos de maclação que ampliam tenacidade mecânica sob frio extremo. \\
\textbf{Revestimentos Térmicos} & Erosão por partículas, atrito de plasma e oxidação extrema. & Formação de filmes superficiais contínuos e passivantes de alumina/sílica. \\
\bottomrule
\end{tabularx}
\legend{Fonte: Elaborado pelos autores (2026), com base em dados compilados da literatura.}
\end{table}

A transposição dessas ligas para prontidão tecnológica de voo enfrenta três desafios principais de engenharia. O primeiro reside na elevada densidade de formulações pioneiras que contêm teores maciços de tungstênio e tântalo; as pesquisas atuais contornam essa penalidade através de ligas leves balanceadas com titânio, zircônio e alumínio. O segundo gargalo, relativo à oxidação intermediária (\emph{pesting}), tem sido controlado pela inclusão sinérgica de formadores de óxidos vítreos protetores. Por fim, a complexidade de usinagem desses materiais ultraduros impulsionou a adoção da manufatura aditiva por fusão em leito de pó a laser (LPBF), viabilizando geometrias aerodinâmicas intrincadas integradas a dutos de resfriamento regenerativo ativo.
